\scanpage{625}
\mspage{616}

% L1
Now let
\[
\tag{\unsure{20}}
\mu_{\tau'}=\{1,\tau'\}\subset\mathcal{G}^{\wedge}
\]
% L3
then \inserted{\unsure{$\mu_{\tau'}$ normalizes $L_\rho$}},
\deleted{\unsure{this}} for
\(\tau'(\rho)=\rho^{-1}\) whence \(\tau'(L_\rho)=L_\rho\)
hence it normalizes \(Z_\rho\unsure{\equiv|=}\operatorname{Cent}_{\mathcal{M}}(L_\rho)\), hence
also \(Z_\rho^!=Z_\rho\cap\mathcal{M}^!\)
(since \(\mathcal{M}^!\) is invariant \unsure{in} \(\mathcal{M}\))
\(=Z_{\rho,0}^!\times\mu=\mathcal{M}'(-\bar{\jmath})\times\mu\).
% L9
\deleted{Let} Moreover
\(\mu_{\tau'}\cap Z_\rho^!\subset
\deleted{\unsure{\{1\}}}\mu_{\tau'}\cap\mathcal{M}^!=\{1\}\).
Let
\[
\tag{21}
\mathcal{N}_\rho^?=\mu_{\tau'}\cdot Z_\rho^!
=\mu_{\tau'}\cdot
\underbrace{(Z_{\rho,0}^!\times\mu)}_{
\substack{\mathcal{M}(-\bar{\jmath})\times\mu\\
\simeq\Gamma_{\mathbb{Q}}^{\prime\sim}\times\mu}}
\quad\deleted{\unsure{\simeq\mathcal{M}_{\mathbb{Q}}\times\mu|\simeq\mathcal{M}_{0}\times\mu}}.
\]
\hfill ($\frac12$ direct product)

% L13
One should be aware that contrary to what
this notation might suggest,
\(\mu_{\tau'}\) does \uline{not} normalize \(Z_{\rho,0}^!\) --- otherwise one
would have a \unsure{decomposition|arrangement} into a product
\[
\mathcal{N}_\rho^?=\mathcal{N}_{\rho,0}^?\times\mu,
\quad\text{where }\mathcal{N}_{\rho,0}^?
=\deleted{\unsure{Z_{\rho,0}^!\times\mu}}
\mu_{\tau'}\cdot\unsure{Z_{\rho,0}^?|Z_{\rho,0}^!},
\]
% L19
but this is not the case. We have an exact sequence
\[
\tag{22}
\diagram{\begin{tikzcd}[ampersand replacement=\&,column sep=large,row sep=small]
1\arrow[r] \& \mu\arrow[r] \& \mathcal{N}_\rho^?\arrow[r]
\& \gamma_{\rho,\sigma}\arrow[r,no head]\arrow[d,equal] \& 1\\
\& \& \& \Gamma_{\mathbb{Q}}^{\sim} \&
\end{tikzcd}}
\]
% L22
so that \(\mathcal{N}_\rho^?\) is a quasi-section
of multiplicity $2$ of the extension
\(\mathcal{N}_\rho^{*}\longrightarrow\gamma_{\rho,\sigma}\), --- of
\(\mathcal{M}\longrightarrow\gamma_{\rho,\sigma}=\Gamma_{\mathbb{Q}}^{\sim}\).
Thus we find

\margin{624}

\scanpage{626}

% L1
\[
\tag{23}
\left\{\begin{aligned}
\mathcal{M}&=\mathcal{N}_\rho^?\wedge_{\mu_2}\mathcal{G}^{+\wedge}\\
\unsure{\mathcal{N}_\rho^*}&=\mathcal{N}_\rho^?\wedge_{\mu_2}
\deleted{\unsure{\mathbb{Z}}}L_\rho
\simeq\mathcal{N}_\rho^?\times
\underbrace{L_{-\rho}}_{\simeq\mathbb{Z}/3\mathbb{Z}}
\end{aligned}\right.
\]
% L3
The \deleted{sequence}\inserted{extension} (22) of
\(\gamma_{\rho,\sigma}\) by \(\mu\) is
the extension denoted by \(\widetilde{\gamma}_{\rho\sigma}\) in the preceding section,
it is surely \uline{\unsure{not} \deleted{trivial}\inserted{\unsure{split}}}.
It gives rise to a canonical \unsure{homomorphism}
\[
\tag{24}
\diagram{\begin{tikzcd}[ampersand replacement=\&,column sep=large]
\gamma_{\rho,\sigma}\arrow[r,"\sim"]
\&\mathcal{N}_\rho^?/\mu\arrow[r]
\&\mathcal{N}_\rho^*/\mu.
\end{tikzcd}}
\]

% L9
\uline{Remark} To do it properly, one would have to
make explicit the operation of \(\tau'\) on
\(Z_\rho^!\simeq Z_{\rho,0}^!\times\mu
\xrightarrow{\sim}\Gamma_{\mathbb{Q}}^{\prime\sim}\times\mu\), in terms of
the operation of the image \(\tau_{\Gamma\deleted{\unsure{\mathbb Q}}}\) of \(\tau'\) on
\(\Gamma_{\mathbb Q}^{\prime\sim}\) by inner \unsure{automorphism}
(N.B. \(\tau_\Gamma\) is none other than the element of
\(\Gamma_{\mathbb Q}\subset\Gamma_{\mathbb Q}^{\sim}\), ``complex conjugation'')
and of a \unsure{homomorphism} from \(\Gamma_{\mathbb Q}^{\prime\sim}\) to \(\mu\)
\inserted{i.e. of a ``quadratic character'' of \(\Gamma_{\mathbb Q}^{\prime\sim}\)}
--- it is a good bet that
this is precisely \(\varepsilon_2\) --- so that
\(\Gamma_{\mathbb Q}^{0\sim}\simeq
\Gamma_{\mathbb Q}^{\prime\sim}\cap\Gamma_{\mathbb Q}^{\prime\prime\sim}\)
appears as invariant under \(\tau'\) \ldots
In short, it would be well to take up again the study
of the extension made in the preceding section,
by \unsure{starting there} from the point of view of the \unsure{associated} \(\alpha,\beta\),

\margin{625}

\scanpage{627}
\mspage{617}

% L1
in the present perspective, probably
simpler. It would be well at the same time
to \deleted{define}\inserted{make explicit} another partial section,
which \deleted{would deserve} would be over
\(\deleted{\unsure{\Gamma}}\Gamma_{\mathbb Q}^{\prime\prime\prime\sim}\),
\[
\Gamma_{\mathbb Q}^{\prime\prime\prime\sim}
=\{x\in\Gamma_{\mathbb Q}^{\sim}\mid
\varepsilon_2(x)\varepsilon_3(x)=1\},
\]
% L6
and which would perhaps deserve the notation
\(\mathcal{M}(\sqrt3)\) \unsure{or at least}
\(\mathcal{M}(q)\), for
a suitable generator \(\unsure{q|\rho}\) of \(\mathbb{Q}(\sqrt3)\) ???
\margin{?\quad\uline{\unsure{Ponder these}}
\uline{\unsure{coincidences.} \unsure{Attentive.|attentively}}}

% L9
I realize that I am
\unsure{talking crap} --- having somewhat lost touch
with the geometric content of my
calculations. One must keep in mind that
\[
\tag{25}
\diagram{\begin{tikzcd}[ampersand replacement=\&]
\mathcal{M}_0^!\arrow[r,"\sim"] \& \mathcal{M}_{0,3}^{!\sim}
\end{tikzcd}}
\]
% L14
corresponds \unsure{up to} $\sim$, to the \(\pi_1\) of \(U_{0,3\mathbb Q}\),
the sections of \(\mathcal{M}_0^!\longrightarrow
\Gamma_{\mathbb Q}^{\sim}\simeq\gamma_{\rho,\sigma}\) thus correspond
(up to $\sim$) to those of
\(\pi_1(U_{0,3\mathbb Q})\) over \(\Gamma_{\mathbb Q}\),
thus (``\unsure{invariably|inevitably|unavowably}'') to the
\inserted{\unsure{rational/$\mathbb Q$}} points of \(U_{0,3\mathbb Q}\)
--- when they are partial sections, to the
points of \(U_{0,3\mathbb Q}\) with values in finite
extensions of \(\mathbb Q\). Thus, \deleted{the} the section
\deleted{\unsure{privileged}}\inserted{considered}
in b) does indeed correspond to the \(\mathrm{pt}\ \frac12\) of \(U_{0,3\mathbb Q}\),
the \deleted{\unsure{quasi-}}section
\inserted{partial \(\mathcal M(-\bar\jmath)\)} considered here in c) corresponds to the
point \(-\bar\jmath\) of \(U_{0,3\mathbb Q}\), which is in
\(U_{0,3\deleted{\unsure{\mathbb Q}}}(\mathbb Q[\sqrt[3]{1}])\),
so the notation \inserted{\(\mathcal M(-\bar\jmath)\)} is reasonable.
As for \(\mathcal N_\rho^{\sim}\),
it is \uline{not} \(\subset\mathcal M_0^!\) --- even on dividing by \(\mu\);
one finds after division by \(\mu\) a section

\margin{626}

\scanpage{628}

% L1
of \(\mathcal M_{0,3}^{\sim}\) over \(\Gamma_{\mathbb Q}^{\sim}\), but this is not a
section of \(\mathcal M_{0,3}^{!\sim}\). Now up to $\sim$, \(\mathcal M_{0,3}^{\sim}\)
is the \(\pi_1\) of \(M_{1,1\mathbb Q}'\), moduli scheme over \(\mathbb Q\)
of elliptic curves \uline{mod symmetry}.
Thus giving this section corresponds to
giving such an ``elliptic curve
mod symmetry'' over \(\mathbb Q\). Passing to the \deleted{sub}
subgroup \(\Gamma_{\mathbb Q}^{\prime\sim}\) --- which corresponds to the extension
of the base field from \(\mathbb Q\) to \(\mathbb Q(j)=\mathbb Q(\sqrt[3]{1})\)
one recovers the image in \(\mathcal M_{0,3}^{!\sim}\) of the preceding \deleted{quasi-}
section \inserted{partial} --- \deleted{it} to \(\mathcal M(-\bar\jmath)\)
--- this thus shows that it is still the ``same'' elliptic
curve. I thus have the impression \deleted{that} however
that over the whole of \(\Gamma_{\mathbb Q}^{\sim}\), this section of
\(\mathcal M_{0,3}^{\sim}\) cannot be lifted to a section
of \(\unsure{\mathcal M_\mu|\mathcal M}\) over \(\Gamma_{\mathbb Q}^{\sim}\)
--- i.e. precisely that the extension \(\mathcal N_\rho^{\unsure{?|*}}\)
of \(\Gamma_{\mathbb Q}^{\sim}\) by \(\mu\) is nontrivial.
Would this \uline{mean} that there is no elliptic curve
defined/\(\mathbb Q\) with invariant $0$? (Something that seems absurd
--- one must be able over \uline{a field} to give any
invariant\ldots). \margin{?}
To be precise, we have the
diagram of moduli schemes over \(\mathbb Q\)
\[
\tag{26}
\diagram{\begin{tikzcd}[ampersand replacement=\&,column sep=large,row sep=large]
\begin{gathered}\text{moduli scheme}\\\text{of elliptic}\\\text{curves}\end{gathered}
\&M_{1,1}\arrow[d,"\substack{\text{gerbe}\\\text{of group }\mu}"]\arrow[dr]
\& \&\\
\begin{gathered}\text{moduli scheme of}\\\text{elliptic curves}\\\text{mod symmetry}\end{gathered}
\&M_{1,1}'\arrow[r]
\&\underset{\unsure{0|\circ}}{\mathbb E^1}
\&\begin{gathered}\text{coarse moduli}\\\text{scheme}\end{gathered}
\end{tikzcd}},
\]
% L26
the fact that a point of \(M_{1,1}'(\mathbb Q)\) does not li-

\margin{627}

\scanpage{629}
\mspage{618}

% L1
ft to a point of \(M_{1,1}\), does not \uline{mean}
that its image in \(\mathbb E^1\) does not lift
to \(M_{1,1}\). To be precise, the \deleted{\illegible{1}}\inserted{diagram}
identifies with:
\[
\tag{27}
\diagram{\begin{tikzcd}[ampersand replacement=\&,column sep=large,row sep=large]
M_{1,1}\simeq(U_{0,3},\widetilde{\mathcal G}_3)
  \arrow[d]\arrow[dr]\&\\
M_{1,1}'\simeq(U_{0,3},\mathcal G_3)\arrow[r]
  \&U_{0,3}/\mathcal G_3\simeq U_{0,3}/\widetilde{\mathcal G}_3
     \simeq\mathbb E^1
\end{tikzcd}},
\]
% L7
where \(\widetilde{\mathcal G}_3\) is the extension of
\(\mathcal G_3\)\inserted{\(=\mathrm{SL}(2,\mathbb Z/2\mathbb Z)\)}
by \(\mu\), defined by
\[
\tag{28}
\begin{aligned}
\widetilde{\mathcal G}_3
 &=\mathcal G^+/\pi_0
 \simeq\mathrm{SL}(2,\mathbb Z)/
 \begin{gathered}
 \text{invariant subgroup}\\
 \text{generated by}\\[-2pt]
 -\ell_0
 \end{gathered}
 \\
 &\simeq\mathrm{SL}(2,\mathbb Z/4\mathbb Z)/
 \begin{gathered}
 \text{invariant}\\
 \text{subgroup}\\
 \text{(of order 4) generated}\\
 \text{by }-\ell_0,\ \text{i.e. \unsure{subgroup|subgroup}}\\
 \text{generated by}\\
 -\ell_0,-\ell_1.
 \end{gathered}
\end{aligned}
\]
% L13
by letting \(\widetilde{\mathcal G}_3\) operate on \(U_{0,3}\) through
\(\mathcal G_3\). A point of \(M_{1,1}'\)\inserted{\unsure{rational/\(\mathbb Q\)}}
identifies with \(\mathcal G_3\)-torsor \inserted{\(T\)} over \(\mathbb Q\), equipped
with a map
\[
\tag{29}
\diagram{\begin{tikzcd}[ampersand replacement=\&]T\arrow[r]\&U_{0,3}\end{tikzcd}}
\]
\margin{\(T\) is the torsor
\inserted{\unsure{under \(\mathrm{SL}(2,\mathbb Z/2\mathbb Z)\)}}
of \unsure{rigidifications|privilegiations} of the points of order $2$ and
\(\widetilde T\) that of ``\unsure{pri-rigidifications|privilegiations|rigidifications}
of the points of order \(\unsure{\leq4|4}\)
\unsure{made precise}'' (the first is defined
for an \unsure{elliptic} curve mod symmetry, and
the $2^{\text{nd}}$ only for a genuine elliptic
curve)}
% L18
compatible with the actions of \(\mathcal G_3\), a
point of \deleted{\(U_{0,3}\)} \(M_{1,1}\) identifies with a \(\widetilde{\mathcal G}_3\)-torsor
\(\widetilde T\) over \(\mathbb Q\), equipped with a map
\[
\tag{30}
\diagram{\begin{tikzcd}[ampersand replacement=\&]\widetilde T\arrow[r,no head]\&U_{0,3}\end{tikzcd}}
\]
compatible with the operation of \(\widetilde{\mathcal G}_3\). \uline{Denoting}
by \deleted{\unsure{\(T\)}}
\[
T=\widetilde T/\mu
\]
the corresponding \(\mathcal G_3\)-torsor, giving
(30) is equivalent to giving (29). Thus the

\margin{628}

\scanpage{630}

% L1
liftings to \(M_{1,1}\) of a point of \(M_{1,1}'\) rational/\(\mathbb Q\)
identify with the \inserted{\(\widetilde{\mathcal G}_3\)-}torsors \(\widetilde T\)
which lift the \(\mathcal G_3\)-torsor \(T\), i.e. (up to isomorphism) they correspond
to the homomorphisms \(\Gamma_{\mathbb Q}\longrightarrow\widetilde{\mathcal G}_3\)
which lift the homomorphism
\(\Gamma_{\mathbb Q}\longrightarrow\mathcal G_3\) which defines \(T\):
\margin{we suppose \unsure{now} chosen
a \deleted{\unsure{geometric}} point of \(T(\overline{\mathbb Q}_0)\)}
\[
\diagram{\begin{tikzcd}[ampersand replacement=\&,column sep=large,row sep=small]
\&\widetilde{\mathcal G}_3\arrow[d]\\
\Gamma_{\mathbb Q}\arrow[r]\arrow[ur,dashed]\&\mathcal G_3
\end{tikzcd}}
\quad -
\]
% L8
the obstruction is, as it should be, an
extension of \(\Gamma_{\mathbb Q}\) by \(\mu_2\), i.e. an \unsure{element}
of \(H^2(\mathbb Q,\mu_2)\simeq {}_2\mathrm{Br}(\mathbb Q)\).
In the case at hand, the image of \(T\) in \(U_{0,3}\)
is made up of \(\{-\bar\jmath,\unsure{-j|-\bar\jmath}\}\)
which is connected over \(\mathbb Q\),
\uline{hence \(\Gamma_{\mathbb Q}\longrightarrow\mathcal G_3\) is surjective}.
\margin{\unsure{corrected|erroneous}, cf. correction below}
In general, \deleted{\unsure{one must}} in the Galois-theoretic translation
in terms of splittings
of \(\mathcal M_{0,3}\) over \(\Gamma_{\mathbb Q}\), the homomorphism
\(\Gamma_{\mathbb Q}\longrightarrow\mathcal G_3\)
is the composite
\[
\diagram{\begin{tikzcd}[ampersand replacement=\&]
\Gamma_{\mathbb Q}\arrow[r,"\unsure{k|\kappa}"]
\&\mathcal M_{0,3}\arrow[r]\&\mathcal G_3
\end{tikzcd}},
\]
% L22
which is trivial iff \deleted{\unsure{\(\Gamma_{\mathbb Q}\)}}
\inserted{\(\unsure{k|\kappa}\)} is a section of
\deleted{of}
\[
\mathcal M_{0,3}^!=\operatorname{Ker}
 (\mathcal M_{0,3}\longrightarrow\mathcal G_3)
\]
over \(\Gamma_{\mathbb Q}\). Here we have a section of
\(\mathcal M_{0,3}^!\) \deleted{over}\inserted{above}
the subgroup \(\Gamma_{\mathbb Q}'\) of index $2$,
\deleted{hence the image of \(\Gamma_{\mathbb Q}'\)}
\(\unsure{k|\kappa}\mid\Gamma_{\mathbb Q}'=1\), hence
\(\unsure{k|\kappa}(\Gamma_{\mathbb Q})\subset\mathcal G_3\)
is of order $2$, we see that

\margin{629}